\section{All finite common-diameter maximizers}\label{sec:common}
Scale the common diameter to one and write
\[
 \mu=\sum_i p_i,\quad V=\sum_i p_i(1-p_i),\quad
 R=\sum_i\{p_i-3p_i^2+4p_i^3-2p_i^4\},\quad
 F_k=\Pp\{\textstyle\sum_i B_i\le k\}.
\]
Since $R\ge V/2$, the ratio is at most $2\sqrt V$.
Its extension by zero at $V=0$ is continuous on the compact probability
cube. Thus $C=C_n^{\mathrm{com}}$ is attained.
At any selected closed integer threshold, define
\begin{equation}\label{eq:deficit}
 H_C(p)=\frac{CR}{V^{3/2}}+
             \Phi((k-\mu)/\sqrt V)-F_k.
\end{equation}
The global definition of $C$ gives $H_C\ge0$ everywhere with $V>0$,
and equality at every maximizing array at its selected threshold.

The probability replacements below have classical precedents.
Hoeffding's pair identity preserves the sum of two probabilities
\cite[equation~(10), pp.~714--716]{Hoeffding1956}.
Mattner's few-moment extremum theorem and next-power-sum argument
\cite[Theorems~1.1--1.2 and Lemma~2.1]{Mattner2022} use related fibers
for separately affine symmetric functions. The present deficit contains
a cubic-moment denominator. We establish its required affinity directly.

\subsection{A uniform third-difference estimate}
For the mass function $c_j$ of a Poisson-binomial sum, extended by zero,
put $\Delta^3c_j=c_j-3c_{j-1}+3c_{j-2}-c_{j-3}$.

\begin{lemma}\label{lem:third-difference}
If that sum has variance $W$, then
\[
                \sup_j|\Delta^3c_j|<\frac{16/5}{(1+W)^{3/2}}.
\]
\end{lemma}
\begin{proof}
For $W\le1/4$, choose the more likely outcome of each bit as a base
configuration. Let $J$ count departures, $\epsilon_i=\min(p_i,1-p_i)$,
$E=\sum\epsilon_i$, and $Q=\sum\epsilon_i/(1-\epsilon_i)$.
Then $E\le2W\le1/2$, $W\le Q\le1$,
\[
 \Pp(J=0)\ge e^{-Q},\qquad
 \Pp(J=1)=Q\Pp(J=0)\ge We^{-W}.
\]
At the two indices where the base state contributes magnitude three,
the one-departure terms have the opposite sign and magnitude at least
one. The base probability is at least $1/2$, so the sign cannot reverse.
Thus the absolute third difference is at most $3-4We^{-W}$ there.
At all other indices it is at most $1+4W$.
For the central bound use $e^{-W}\ge1-W$ and
\[
 (1+W)^{-3/2}\ge h(W):=
       1-3W/2+15W^2/8-35W^3/16.
\]
The polynomial $(16/5)h(W)-(3-4W+4W^2)$ is decreasing on $[0,1/4]$:
its derivative is bounded above by
$-1/2+13W/4-315W^2/16<0$.
At $W=1/4$ it equals $1/64$.
The outer estimate is at most two, below $(16/5)(5/4)^{-3/2}$.

For $W\ge1/4$, Fourier inversion gives
\[
 |\Delta^3c_j|\le I(W):=
 \frac{16}{\pi}\int_0^{\pi/2}\sin^3u\,e^{-2W\sin^2u}\,du.
\]
Write $J_r=\int_0^{\pi/2}\sin^ru\,e^{-2W\sin^2u}\,du$.
Integration of $(\sin^4u\cos u\,e^{-2W\sin^2u})'$ gives
$(5+4W)J_5=4J_3+4WJ_7$. Consequently
\[
 \frac{d}{dW}\log\{I(W)(1+W)^{3/2}\}
 \le-\frac8{5+4W}+\frac3{2(1+W)}<0.
\]
At $W=1/4$, the exponential upper Taylor polynomial through degree six
and the exact odd sine integrals give
\[
 I(1/4)(5/4)^{3/2}
 \le\frac{10\sqrt5}{\pi}\,
       \frac{911264}{2027025}<\frac{16}{5}.
\]
The last strict inequality follows by squaring and using
$\pi>6283/2000$; it is one of the rational scalar controls in
the computational supplement.
\end{proof}

\subsection{Reducing the number of distinct probabilities}
Fix a maximizing array and its mean, variance and closed threshold.
Consider the compact zero set
\[
 Z=\{p\in[0,1]^n:\ \sum p_i=\mu,\
                  \sum p_i(1-p_i)=V,\ H_C(p)=0\}.
\]
Let $K$ be its connected component containing the array, and maximize
$\sum p_i^3$ on $K$. For three selected probabilities, let $e_1,e_2,e_3$
be their elementary symmetric functions. Their generating polynomial is
\[
 1+e_1(w-1)+e_2(w-1)^2+e_3(w-1)^3.
\]
Holding $e_1,e_2$ fixed preserves their mean and variance. Their
contribution to $R$ is
\begin{align*}
 R_3={}&e_1-3e_1^2+4e_1^3-2e_1^4
 +(6-12e_1+8e_1^2)e_2-4e_2^2+4(3-2e_1)e_3.
\end{align*}
Both $F_k$ and $R$ are affine in $e_3$ on this fiber. If its three
roots are distinct and interior, $e_3$ can move in both directions.
A nonzero slope contradicts $H_C\ge0$. A zero slope gives a continuous
path in $Z$, hence in $K$, on which $\sum p_i^3=e_1^3-3e_1e_2+3e_3$
increases. The chosen point therefore has at most two distinct active
probabilities.

Suppose they are $a<b$, each occurring at least twice. The triples
$(a,a,b)$ and $(a,b,b)$ are opposite endpoints of their three-root
fibers. If $d_j$ is the mass function of the other coordinates, the
fiber slope is
\[
 L=\frac{4C(3-2e_1)}{V^{3/2}}+\Delta^2d_k.
\]
The two endpoint conditions give $L_{aab}\le0\le L_{abb}$.
Delete two bits of each type, let $c_j$ be the remaining mass function,
and let $W$ be its variance. Direct subtraction gives
\[
 L_{abb}-L_{aab}
       =(b-a)\left\{\Delta^3c_k-\frac{8C}{V^{3/2}}\right\}.
\]
Since $V\le W+1$, Lemma~\ref{lem:third-difference} contradicts these
signs whenever $C\ge2/5$. Thus one of the two multiplicities is one.
At $n=3$ this is automatic. At $n=4$, the only remaining repeated
pattern is $(2,2)$; its deleted law is deterministic, so
$|\Delta^3c_k|\le3$, whereas $8C/V^{3/2}\ge8C>3$ when
$C\ge199/500$.

\subsection{Strict curvature of the remaining fiber}
The remaining unequal pattern is $\Bin(M,p)+\Bern(x)$, $p\ne x$.
We record its exact curvature and the finite inequality used to exclude it.
Let $q=1-p$, $m=M-1$, $d_j=\Pp(\Bin(m,p)=j)$, $W=mpq$, and
\[
 V=x(1-x)+Mpq,\quad z=\frac{k-x-Mp}{\sqrt V},\quad
 \xi=d_{k-1}-d_k,\quad
 S_2=(1-2x)^2+(1-2x)(1-2p)+(1-2p)^2.
\]
Preserve the sum of the exceptional mass and one mass $p$, write them
as $\bar p+d,\bar p-d$, and set $u=d^2$. Along this path,
\[
 F_k'=\xi,\qquad V'=-2,\qquad R'=-2S_2,\qquad R''=-8.
\]
At stationarity of $H_C$, differentiating its smooth terms and
eliminating $R$ gives
\begin{equation}\label{eq:fiber-curvature}
 V^2H_C''=5V\xi-C\left(8\sqrt V+\frac{2S_2}{\sqrt V}\right)
                       -z\phi(z)(2+z^2).
\end{equation}
A minimum requires this expression to be nonnegative.

\begin{lemma}[The certified curvature inequality]\label{lem:certified-curvature}
Put $c_*=199/500$. For
\[
 0<p,x<1,\quad pq\le W\le3,\quad0\le f\le1,\quad
 A_0=W-(1-f)q^2\ge0,
\]
define
\[
 B_0=p^2+(1-f)(W+2pq),\quad V=x(1-x)+W+pq,
 \quad z=\frac{f(W+pq)/(1-qf)-x}{\sqrt V}.
\]
Then
\begin{equation}\label{eq:compact-curvature}
 \frac{2c_*}{\sqrt V}\{4W+2+(1-2x)(1-2p)\}
 +z\phi(z)(2+z^2)
 -\frac{5VfB_0}{(2-f)B_0+A_0}>0.
\end{equation}
Also every two-bit unit-diameter array has ratio less than $c_*$.
\end{lemma}

The lemma is a finite arithmetic premise. Its compact interval cover
at 192 Arb bits \cite{Johansson2017} processed 550563 boxes for
\eqref{eq:compact-curvature} and 13641 boxes for the two-bit assertion,
with no unresolved box. The archived records bind the evaluated
formulas and the declared domains. Below we explain the covering
charts and the analytic complements; rebuilding those interval covers
is distinct from checking the saved hashes and coverage metadata.

For $V\le1/5$, the first term in \eqref{eq:compact-curvature} is at
least $8c_*/\sqrt3$, the normal term is at least $-4/(e\sqrt\pi)$,
and the final term is at most one. Rational controls give
\begin{equation}\label{eq:scalar-margins}
 \frac{8c_*}{\sqrt3}-\frac4{e\sqrt\pi}-1>\frac1{247},\quad
 8c_*\sqrt3-\frac{10}{\pi}-\frac4{e\sqrt\pi}>1,\quad
 8c_*\sqrt{7/2}-\frac{35}{3\pi}-\frac4{e\sqrt\pi}>\frac{44}{45}.
\end{equation}
For $V\ge1/5$ use
$z\phi(z)(2+z^2)\ge\sqrt{2/\pi}\min(z,0)$.
The following maps from $[0,1]^4$, with fourth variable $x$, cover
the entire remaining domain:
\begin{align*}
 &p=(1+s_p)/2,\quad W=pq+s_W(3-pq),\quad f=s_f;\\
 &p=s_p/2,\quad W=q^2+s_W(3-q^2),\quad f=s_f;\\
 &p=s_p/2,\quad W=pq+s_W(q^2-pq),\quad
                    f=1-(1-s_f)W/q^2.
\end{align*}
The first covers $p\ge1/2$. The last two split $p\le1/2$ at $W=q^2$;
below that interface the last map imposes exactly $f\ge1-W/q^2$.
Interval division is used only with a strictly positive denominator.
Otherwise the universal neighbor-fraction bound one, or $y\ge0$ at a
singular endpoint, supplies a valid lower bound. A box is accepted only
when its outward lower endpoint is positive. The two-bit cover includes
the entire probability square and all three closed integer thresholds.
The zero-variance corners use $\mathcal C\le2\sqrt V$ before division.

To apply the lemma to \eqref{eq:fiber-curvature}, first suppose $M\ge2$.
If $\xi\le0$, then $\max(V,S_2)\ge1/3$ and
\[
 C(8\sqrt V+2S_2/\sqrt V)+z\phi(z)(2+z^2)
 \ge8c_*/\sqrt3-4/(e\sqrt\pi)>0.
\]
Thus $H_C''<0$. If $\xi>0$, binomial recurrence gives
\[
 y=k-(m+1)p>0,\quad f=y/(kq)\in(0,1],\quad
 d_k=(1-f)d_{k-1},\quad d_{k-2}=r_0d_{k-1},
\]
where
\[
 r_0=A_0/B_0,\qquad y=f(W+pq)/(1-qf).
\]
The sum of these three masses is at most one, so
$\xi\le f/(2-f+r_0)$. For $W\le3$, Lemma~\ref{lem:certified-curvature}
therefore makes \eqref{eq:fiber-curvature} negative. The coefficient
of $C$ is positive, so the same conclusion holds for every actual extremal ratio
$C\ge c_*$.
For $W\ge3$, Fourier inversion gives
\[
 |\xi|\le\frac4\pi\int_0^{\pi/2}\sin u\,e^{-2W\sin^2u}\,du
                        <\frac2{\pi W}.
\]
The negative of the right side of \eqref{eq:fiber-curvature} is at least
$8c_*\sqrt V-4/(e\sqrt\pi)-10V/(\pi W)$.
It is concave in $V\in[W,W+1/2]$. Both endpoint values increase with
$W$, and \eqref{eq:scalar-margins} gives positive bounds at $W=3$.
For $M=1$ the two-bit assertion of the lemma excludes the pattern at
price $C\ge c_*$.

\subsection{From one equal-mass point to every maximizer}
The cubic maximizer on $K$ is now an IID group plus deterministic bits.
The number $m$ of active bits is at least three, by
Lemma~\ref{lem:certified-curvature}. We prove this point is isolated in $Z$.
For a pair in the IID group, use $p+d,p-d$ and $u=d^2$.
Nonnegativity gives $\Lambda:=\partial_uH_C(0)\ge0$.
If $\Lambda=0$, \eqref{eq:fiber-curvature} applies at the diagonal
with deleted law $\Bin(m-2,p)$. The compact lemma or its analytic
complement makes $\partial_u^2H_C(0)<0$, contradicting $H_C\ge0$
for small positive $u$. Hence $\Lambda>0$.
Permutation symmetry makes the exchange quadratic form
$(\Lambda/2)\|y_{\mathrm{cen}}\|^2$ on the mean-zero directions
of the active probabilities.

For a nearby point in $Z$, let $a_i\ge0$ activate zero probabilities,
$b_j\ge0$ deactivate probability-one coordinates, and $y_l$ change
the $m$ active probabilities. Put $a=\sum a_i$, $b=\sum b_j$, $t=a+b$.
The two constraints give exactly
\[
 \sum_l y_l=b-a,\qquad
 \sum_l y_l^2=2pa+2(1-p)b-\sum_i a_i^2-\sum_j b_j^2.
\]
Thus $\|y_{\mathrm{cen}}\|^2\ge c_pt-O(t^2)$, with $c_p>0$.
The endpoint linear coefficients of $H_C$ are nonnegative: activate
one endpoint and compensate uniformly in the active group to preserve
the mean, then apply $H_C\ge0$ on that feasible path.
The Taylor expansion on the fixed mean/variance fiber consequently has
nonnegative endpoint linear terms, the positive exchange term, and
remaining terms $O(t^{3/2})$. The uniform active direction is $O(t)$,
the exchange direction is $O(\sqrt t)$, so these remainder orders also
include their mixed terms. For small $t>0$ this gives $H_C>0$.
If $t=0$, the two constraints force $y=0$. The IID point is isolated.
Its connected component $K$ is therefore a singleton, and the original
maximizer was the same IID group plus zeros.

The required lower prices follow from two exact admissible examples:
\[
 \begin{array}{c|c|c|c}
 \text{active count}&p&\text{closed raw threshold}&\text{ratio lower bound}\\\hline
 3&43/100&1&0.3981444741>199/500\\
 5&9/20&2&0.4009890598>2/5.
 \end{array}
\]
Their binomial probabilities are rational. Rational square-root bounds,
Machin bounds for $\pi$, and the integrated Gaussian Taylor polynomial
$x-x^3/6+x^5/40-x^7/336$ give the displayed strict lower bounds.
Padding preserves them. The three-bit example handles $n=3,4$;
the five-bit example supplies the $2/5$ hypothesis for the repeated-type
argument at every $n\ge5$.

\subsection{The remaining two-coordinate case}
For two probabilities $p,x$, set $s=p+x$, $w=px$ and
\[
 V=s-s^2+2w,\quad
 R=s-3s^2+4s^3-2s^4+(6-12s+8s^2)w-4w^2.
\]
At rank $k=0,1,2$, the closed CDFs are $1-s+w$, $1-w$, $1$.
Let $z=(k-s)/\sqrt V$, $D=F_k-\Phi(z)$ and $C=DV^{3/2}/R$.
Off the diagonal the map $(p,x)\mapsto(s,w)$ is invertible, so an
interior local maximum requires both derivatives to vanish.
With $F_w=(1,-1,0)_k$, $F_s=(-1,0,0)_k$,
$V_s=1-2s$, $R_w=6-12s+8s^2-8w$ and
$R_s=(1-2s)^3+(-12+16s)w$, their positive-denominator numerators are
\begin{align*}
 G_w={}&R(VF_w+z\phi(z)+3D)-DVR_w,\\
 G_s={}&2VRF_s+2\sqrt V R\phi(z)+Rz\phi(z)V_s
                           +3DRV_s-2DVR_s.
\end{align*}
At $G_w=0$,
\[
 \frac{C_{ww}}C=\frac{N}{DV^2R},\quad
 N=R\{6F_wV+z\phi(z)(3+z^2)+3D\}+8DV^2.
\]
The finite critical certificate covers each owner square by
$C<37/100$, $0\notin G_w$, $0\notin G_s$, or $N>0$.
A separate diagonal cover uses only the ratio bound, since the
$(s,w)$ chart is singular there. These covers processed respectively
4683 and 195 boxes, with no unresolved box. They exclude every
two-active-bit common-diameter maximizer at ratio at least $37/100$.
Boundary derivatives are not assumed to vanish; the cover bounds their
limits only to enclose the open stationary domain.
The exact one-bit example $p=2/5$ gives a ratio greater than $37/100$.
Compact attainment therefore leaves exactly one active bit at $n=2$.
The $n=1$ case is immediate. This completes
Theorem~\ref{thm:common}.
